\section{Preparation and experimental design}\label{sec:design}

\subsection{Preparation: literature, derivations and a 2D pilot study}\label{sec:prep}
\paragraph{Reading and derivations.} One author worked through the flow-matching papers \citep{lipman2023fm,liu2023rf,albergo2023si,tong2024otcfm}, the other the diffusion papers \citep{song2019ncsn,ho2020ddpm,song2021sde,song2021ddim,luo2022understanding}, each re-deriving the main formulas in a common notation (\texttt{docs/notes/}). Every derivation in Section~\ref{sec:theory} was written out step by step (\texttt{docs/notes/derivations.md}) and, where possible, checked against an independent computation: the DDIM--Euler identity on a 1D mixture, the Gelbrich bound on the MNIST pixel covariance, the $\sqrt\kappa$ grid against the exact field of two 1D targets, and the TV noise floor against a multinomial simulation. We also collected published CIFAR-10 numbers with their exact settings (Table~\ref{tab:literature}).

\paragraph{2D pilot study.} We first trained I-CFM, OT-CFM and DDPM on four standardised 2D datasets (8 Gaussians, checkerboard, two moons, Swiss roll) with an MLP (three hidden layers of width 256, SiLU, 64-dimensional sinusoidal time embedding; Adam, learning rate $10^{-3}$, 20,000 iterations at batch 256). Samples from 2,000 noise vectors are compared with 2,000 data points by the exact 2-Wasserstein distance $W_2$ (POT, \citealp{flamary2021pot}); two independent data samples are 0.04--0.15 apart. With one Euler step OT-CFM reaches $W_2=0.12$--$0.21$, I-CFM 1.28--1.37 and DDIM 7.2--12.6; at 100 steps they reach 0.09--0.20, 0.10--0.31 and 0.17--0.47, against 0.15--0.29 for DDPM with 1,000 stochastic steps. Figure~\ref{fig:2d} shows the mechanism of Proposition~\ref{prop:onestep}. % results/e1/summary.csv
\begin{figure}[t]
\centering
\includegraphics[width=0.62\linewidth]{figs/fig_2d_paths.pdf}
\caption{2D pilot study on the two-moons data (seed 0): 100-step Euler trajectories from Gaussian noise (dots); crosses mark where a single Euler step lands. I-CFM's one-step samples collapse to the data mean (Proposition~\ref{prop:onestep}); OT-CFM's straight paths make one step nearly exact. Minibatch OT saves 87--93\% of the pairing cost in 2D, 4.5\% on MNIST.}
\label{fig:2d}
\end{figure}

\paragraph{From the pilot to the research questions.} On MNIST (shown in our class talk) OT-CFM was better than I-CFM only at one to five steps. Three observations shaped the final design. (1) Proposition~\ref{prop:bound} explains the gap: pairing can save about the whole cost in 2D but at most 26\% on MNIST, so work on the coupling alone could not be expected to help much. (2) Two DDIM seeds with almost identical training loss (0.0215 and 0.0213) had FID 12.2 and 32.8 at 100 steps, and the better one had the worse class balance (Section~\ref{sec:c1}), so the evaluation error had to be measured before comparing improvements. (3) Inference-side changes are free to try on an existing model, training-side changes cost a training run, so both must be compared with their cost. We dropped a CIFAR-10 study (too expensive to train several variants), EMA weights and learning-rate schedules (they change every model and do not target the number of steps), and consistency or MeanFlow models (different objectives).

\subsection{Data, network, training and samplers}\label{sec:setup}
\paragraph{Models.} MNIST \citep{lecun1998mnist}: 60,000 training images of $28\times28$ pixels scaled to $[-1,1]$, no augmentation. All models use the same U-Net, the ADM architecture of \citet{dhariwal2021adm} as implemented in \texttt{torchcfm} \citep{tong2024otcfm} (32 base channels, one residual block per resolution, 1,075,361 parameters, PyTorch, \citealp{paszke2019pytorch}), trained with Adam \citep{kingma2015adam} at learning rate $10^{-3}$, batch 128, for 50 epochs (468 steps each), without EMA or schedule; the last-epoch weights are evaluated. OT-CFM re-pairs each batch by an exact OT plan for the squared Euclidean cost (POT's network-simplex solver). The diffusion reference is DDPM \citep{ho2020ddpm} with linear $\beta_t$ from $10^{-4}$ to $0.02$, $T=1000$ and the $\epsilon$-loss, sampled with deterministic DDIM ($\eta=0$) on evenly spaced timesteps. I-CFM and OT-CFM have three training seeds, DDPM two.
\paragraph{Training-side variants.} Each changes one setting of the seed-0 model: \emph{OT batch} 512 or 2,048 (the plan is computed on the larger batch, which is then split into batches of 128, so the number of updates is unchanged); \emph{time distribution} logit-normal or the $\kappa$-mixture of Section~\ref{sec:tdist}, sampled by inverting its piecewise-linear CDF; \emph{reflow}: 60,000 pairs $(x_0,\hat x_1)$ from integrating the seed-0 model with dopri5 from fresh noise, then 50 epochs of fine-tuning from the base weights with a new optimiser.
\paragraph{Solvers and grids.} Euler, Heun, RK4, Heun* and RK4* (Section~\ref{sec:t1}); the reference is the adaptive Dormand--Prince 5(4) method \citep{dormand1980rk} from \texttt{torchdiffeq} \citep{chen2018node} with $\mathrm{rtol}=\mathrm{atol}=10^{-5}$ (about 270 NFE). Grids: uniform; the shift family $t'=st/(1+(s-1)t)$ of \citet{esser2024sd3} with $s=2^e$, $e\in\{-3,-2.75,\dots,3\}$; and the $\sqrt\kappa$ grid \eqref{eq:sqrtk}. We estimate $\kappa(t)$ from 256 dopri5 trajectories per model: $v(x(t),t)$ is recorded at 401 equally spaced times, $\|a(t)\|$ is the per-pixel RMS of finite differences, averaged over trajectories, smoothed over neighbouring intervals and floored at $10^{-3}$ of its maximum. As a data-free alternative, a \emph{calibrated} grid chooses for every solver and NFE the shift grid with the smallest endpoint error against dopri5 on 256 paths (noise seed 300), reported on a separate set (seed 301); it tests whether ODE accuracy predicts sample quality.

\subsection{Evaluation protocol}\label{sec:eval}
\paragraph{FID, efficiency and straightness.} FID \citep{heusel2017fid} fits Gaussians $(\mu_r,\Sigma_r)$, $(\mu_g,\Sigma_g)$ to Inception-v3 features \citep{szegedy2016inception} of real and generated images and reports $\|\mu_r-\mu_g\|^2+\operatorname{tr}(\Sigma_r+\Sigma_g-2(\Sigma_r\Sigma_g)^{1/2})$, the squared Wasserstein distance between the Gaussians (Lemma~\ref{lem:gelbrich} with equality). We use clean-fid \citep{parmar2022cleanfid} (\texttt{clean} mode) with 10,000 samples against the 60,000 training images; real-versus-real FID is 0.54. All samplers on a model start from the same noise (seed 100), so their differences are paired. Because the estimate is biased upward roughly as $1/n$ \citep{chong2020fid}, we also fit $\mathrm{FID}(n)=\mathrm{FID}_\infty+c/n$. NFE@20 is the NFE at which FID first reaches 20, interpolated linearly in $\log\nfe$; $R$ (Section~\ref{sec:reflow}) uses 256 paths of 100 Euler steps; training cost is wall time relative to a base model under the same throttling.
\paragraph{Class coverage: the course method.} FID cannot say \emph{which} kinds of digits a model over- or under-produces. Our course method, $k$-means clustering \citep{lloyd1982kmeans} followed by a contingency-table test, answers this without labels. (1) Split the training images at random into a fitting and a reference half. (2) Fit $k$-means ($k=10$, 10 initialisations, scikit-learn, \citealp{pedregosa2011sklearn}) on the fitting half only, so the clusters do not depend on any generative model; another initialisation and $k=20$ are robustness checks. (3) Assign the reference half and 10,000 generated images to the nearest centre and count per cluster. (4) Compare the generated shares $\hat p_k$ with the reference shares $p_k$ ($m=30{,}000$) by total variation and by Pearson's two-sample $\chi^2$ test of homogeneity on the $2\times k$ table ($k-1$ degrees of freedom, \citealp[\S3.2]{agresti2013cda}), with Cohen's $w=\sqrt{\chi^2/n}$ \citep{cohen1988power} and bootstrap intervals \citep{efron1979bootstrap} that resample both count vectors (1,000 resamples). The two-sample form matters: the usual statistic $\sum_k(O_k-np_k)^2/(np_k)$ treats the estimated $p_k$ as known and overstates the evidence. (5) Run the same pipeline on the 10,000 MNIST test images as a perfect generator. (6) For interpretation only, label generated images with an RBF-SVM trained on 10,000 real images (test accuracy 0.970). Even a perfect model has positive TV, because both share vectors are estimated: % results/c1_reliability/kmeans_c1.json
\begin{proposition}[TV noise floor]\label{prop:tvfloor}
Let $\hat p$, $\tilde p$ be the cluster shares of $n$ and $m$ independent draws from the same distribution $p$ over $k$ clusters. As $n,m\to\infty$,
\begin{equation}\label{eq:tv}
\E\,\mathrm{TV}=\E\,\tfrac12\textstyle\sum_k|\hat p_k-\tilde p_k|\;\approx\;\tfrac12\sum_k\sqrt{\tfrac{2}{\pi}\,p_k(1-p_k)\big(\tfrac1n+\tfrac1m\big)} .
\end{equation}
\end{proposition}
\begin{proof}
Each $\hat p_k-\tilde p_k$ has mean 0 and variance $p_k(1-p_k)(1/n+1/m)$ and is asymptotically normal (central limit theorem); for $Z\sim\mathcal N(0,\sigma^2)$, $\E|Z|=\sigma\sqrt{2/\pi}$. Linearity of expectation over $k$ gives the result without independence between clusters.
\end{proof}
With $k=10$, $p_k\approx0.1$, $n=10{,}000$, $m=30{,}000$: $\tfrac12\cdot10\cdot\sqrt{\tfrac2\pi\cdot0.09\cdot(10^{-4}+3.3\times10^{-5})}\approx0.014$; a simulation gives 0.0139, the real test set 0.017.
\paragraph{Statistical decision rule.} We model the FID of training seed $s$ under noise draw $r$ as $F_{s,r}=\mu+a_s+e_{s,r}$ with independent seed effects (variance $\sigma^2_{\text{seed}}$) and noise effects (variance $\sigma^2_{\text{noise}}$), so $\hat\sigma^2_{\text{seed}}=s^2_{\text{between}}-\hat\sigma^2_{\text{noise}}$ (two degrees of freedom from three seeds). A separately trained variant differs from its base also by a new seed effect, so it counts as better or worse only if the difference exceeds $2\hat\sigma_{\text{seed}}$ of its base method at that NFE (the \emph{threshold}). Two samplers on one model and the same noise differ by $D=F^A-F^B$, which contains no seed effect, so they need only exceed the sampling noise and keep their sign on the other seeds. With three seeds the smallest one-sided sign-test $p$-value is $0.5^3=0.125$, so we report sign agreement, not significance. The experiments and this rule were written down (\texttt{docs/FINAL\_REPORT\_PLAN.md}) before the runs; the $t=1$ diagnostic, Heun* and RK4* were added after the first solver results.

\subsection{Compute and reproducibility}\label{sec:compute}
All experiments ran on one consumer GPU, one job at a time, throttled to about half duty cycle and paused above 72\,$^\circ$C (\texttt{code/common/throttle.py}) so that the computer stayed usable; wall times are therefore accurate only to tens of percent (two identical reflow retrainings took 1.08 and 1.56\,h). A base model takes 33--35 minutes unthrottled, 70--72 throttled; training and pair generation took about 20 GPU-hours in total. The computer lost power twice during GPU jobs, so every script skips finished work, training checkpoints every epoch, and \texttt{code/run\_queue.sh} runs the final experiments in order and resumes after an interruption. \texttt{report/make\_figs.py} and \texttt{make\_tables.py} regenerate every figure and table from the result files, and \texttt{code/README.md} gives the command for every experiment. % results/e2/*_train.json; results/c5_reflow/pairs_*_s0.json
